%! Unit 1: Asking Questions with Data: Study Design — Summative Assessment — Practice Test --- Study Copy (KEY)
\documentclass[10pt]{article}
\usepackage{saar-article}
\usepackage{saar-key}   % pulls in -boxes; do NOT also load -boxes

% Work blocks on a test need handwriting room, not typeset spacing. This moves
% the blank and the key together, so it cannot open a page mismatch.
\setlength{\workrowsep}{5pt}

% Part-divider strip
\newcommand{\parthead}[1]{%
  \vspace{6pt}\noindent
  \begin{headlinebox}{cerulean}{\color{white}\bfseries #1}\end{headlinebox}%
  \vspace{4pt}}

% NO teachernote here: the practice test rides in the STUDENT packet, and this
% key must stay the same length as its blank. Answer rationale and Part D
% scoring live on page 2 of unit01/unit_cover_key/main.tex (key packet only).

\begin{document}

\pageheader{Unit 1: Asking Questions with Data: Study Design}{Practice Test --- Study Copy --- Answer Key}
\namedateperiod

\vspace{0.10in}
\begin{remindbox}
\small
\textbf{This is a practice test.} It mirrors the real Unit 1 test in length, format,
and the ideas it covers, but the numbers and contexts differ. Work every problem
with no notes, then check your answers against the key.
\textbf{100 points:} Part A $14$, Part B $16$, Part C $40$, Part D $30$.
\end{remindbox}

% ======================================================================
\parthead{Part A --- Vocabulary \hfill 14 questions, 1 point each}

\small
\textbf{Set 1 --- Samples and Bias.} Write the letter of the matching description
next to each term.

\vspace{2pt}
\begin{minipage}[t]{0.36\linewidth}
\begin{enumerate}[leftmargin=1.9em, itemsep=3.5pt, topsep=1pt]
  \item \ans{C} Population
  \item \ans{F} Parameter
  \item \ans{A} Statistic
  \item \ans{B} Sampling frame
  \item \ans{G} Stratified sample
  \item \ans{D} Cluster sample
  \item \ans{E} Undercoverage
\end{enumerate}
\end{minipage}\hfill
\begin{minipage}[t]{0.61\linewidth}
\begin{enumerate}[label=\Alph*., leftmargin=1.9em, itemsep=3.5pt, topsep=1pt]
  \item A number that describes a sample.
  \item The list of individuals a sample can actually be drawn from.
  \item Every individual the study wants to describe.
  \item A few whole groups are chosen at random, and everyone in them is studied.
  \item Part of the population has no chance of being chosen at all.
  \item A number that describes a population.
  \item The population is split into groups, and some are chosen at random from
        every group.
\end{enumerate}
\end{minipage}

\vspace{6pt}
\textbf{Set 2 --- Studies and Experiments.} Same directions.

\vspace{2pt}
\begin{minipage}[t]{0.36\linewidth}
\begin{enumerate}[leftmargin=1.9em, itemsep=3.5pt, topsep=1pt, start=8]
  \item \ans{E} Observational study
  \item \ans{A} Experiment
  \item \ans{C} Treatment
  \item \ans{G} Control group
  \item \ans{B} Lurking variable
  \item \ans{F} Confounding variable
  \item \ans{D} Nonresponse
\end{enumerate}
\end{minipage}\hfill
\begin{minipage}[t]{0.61\linewidth}
\begin{enumerate}[label=\Alph*., leftmargin=1.9em, itemsep=3.5pt, topsep=1pt]
  \item The researcher assigns individuals to groups and imposes a condition.
  \item A variable outside the study that could be moving both of the variables
        being compared.
  \item The specific condition a group is given.
  \item Individuals were chosen for the sample but never answered.
  \item The researcher measures or surveys without assigning anyone to a group.
  \item The groups inside the experiment differ in a second way too, so you
        cannot tell which difference produced the result.
  \item The baseline group, which gets no treatment or the usual condition.
\end{enumerate}
\end{minipage}
\normalsize

% ======================================================================
\boxguard[8]
\parthead{Part B --- Multiple Choice \hfill 8 questions, 2 points each}

\small
\begin{enumerate}[leftmargin=2.2em, itemsep=5pt, topsep=2pt, start=15]

\item Riverbend High School records each student's \textbf{grade level} as
$9$, $10$, $11$, or $12$. This variable is
\begin{enumerate}[label=(\Alph*), leftmargin=1.8em, itemsep=1pt, topsep=1pt]
  \item quantitative, because the values are numbers.
  \item categorical, because the numbers are labels for groups.
        \textcolor{keyred}{\textbf{$\leftarrow$ correct}}
  \item quantitative, because you can compute a mean of them.
  \item neither, because grade level is not really data.
\end{enumerate}

\item District records show that \textbf{exactly $20\%$ of all $1{,}500$ Harbor
Point members} swim before 9 a.m. That $20\%$ is a
\begin{enumerate}[label=(\Alph*), leftmargin=1.8em, itemsep=1pt, topsep=1pt]
  \item statistic, because it is written as a percent.
  \item statistic, because $1{,}500$ is a large number.
  \item parameter, because it describes every member of the population.
        \textcolor{keyred}{\textbf{$\leftarrow$ correct}}
  \item parameter, because it came out to a round number.
\end{enumerate}

\item Which sentence states the difference between a stratified sample and a
cluster sample?
\begin{enumerate}[label=(\Alph*), leftmargin=1.8em, itemsep=1pt, topsep=1pt]
  \item Stratified takes some from every group; cluster takes everyone from a
        few groups. \textcolor{keyred}{\textbf{$\leftarrow$ correct}}
  \item Stratified takes everyone from a few groups; cluster takes some from
        every group.
  \item Stratified uses chance to choose; cluster does not use chance at all.
  \item There is no difference --- the two names describe the same procedure.
\end{enumerate}

\item A survey is conducted by calling landline telephones. Households that own
no landline can never be reached. This is an example of
\begin{enumerate}[label=(\Alph*), leftmargin=1.8em, itemsep=1pt, topsep=1pt]
  \item nonresponse, because those households did not answer.
  \item undercoverage, because those households could never be chosen.
        \textcolor{keyred}{\textbf{$\leftarrow$ correct}}
  \item sampling variability, because samples differ from each other.
  \item a response bias, because the question was worded badly.
\end{enumerate}

\item A survey's method causes it to overestimate by $16$ percentage points. The
researcher repeats it using \textbf{ten times as many people} and the same
method. The new estimate will most likely
\begin{enumerate}[label=(\Alph*), leftmargin=1.8em, itemsep=1pt, topsep=1pt]
  \item land very close to the true value, because the sample is large.
  \item underestimate, because large samples pull the other way.
  \item be impossible to predict in any direction.
  \item overestimate by about the same $16$ points.
        \textcolor{keyred}{\textbf{$\leftarrow$ correct}}
\end{enumerate}

\item In a random sample of Riverbend students, those who use the Study Center
pass \emph{fewer} classes than those who do not. This result shows
\begin{enumerate}[label=(\Alph*), leftmargin=1.8em, itemsep=1pt, topsep=1pt]
  \item that the Study Center causes students to pass fewer classes.
  \item that passing classes causes students to avoid the Study Center.
  \item an association; a lurking variable could explain the whole thing.
        \textcolor{keyred}{\textbf{$\leftarrow$ correct}}
  \item nothing at all, because the sample was chosen at random.
\end{enumerate}

\item In an experiment, the purpose of the \textbf{control group} is to
\begin{enumerate}[label=(\Alph*), leftmargin=1.8em, itemsep=1pt, topsep=1pt]
  \item provide a baseline the treatment group can be compared against.
        \textcolor{keyred}{\textbf{$\leftarrow$ correct}}
  \item give some subjects a break from the study.
  \item make the total number of subjects larger.
  \item remove the need for random assignment.
\end{enumerate}

\item A researcher wants to know whether smoking causes lung disease in
teenagers. No experiment is run. The best reason is that the experiment would be
\begin{enumerate}[label=(\Alph*), leftmargin=1.8em, itemsep=1pt, topsep=1pt]
  \item impractical, because volunteers are hard to find.
  \item unethical, because it would assign teenagers to smoke.
        \textcolor{keyred}{\textbf{$\leftarrow$ correct}}
  \item impossible, because smoking cannot be measured.
  \item unnecessary, because observational studies prove causation.
\end{enumerate}

\end{enumerate}
\normalsize

% ======================================================================
\boxguard[8]
\parthead{Part C --- Short Answer \& Computation \hfill 8 questions, 5 points each}

\small
\begin{enumerate}[leftmargin=2.2em, itemsep=9pt, topsep=2pt, start=23]

% ---------------------------------------------------------------- C1 (1.1)
\item \textbf{Bayside Middle School} has $750$ students. For each student the
principal records: grade level, minutes spent on homework last night, whether
the student rides the bus, and favorite lunch entr\'ee.

\vspace{2pt}
(a) Who are the \textbf{individuals} in this data set? \ans{The 750 students}

\vspace{3pt}
(b) Classify each variable.

\vspace{2pt}
\renewcommand{\arraystretch}{1.25}
\begin{tabularx}{\linewidth}{@{} X c @{}}
\rowcolor{frost}
\textbf{Variable} & \textbf{Categorical or quantitative?} \\
\midrule
Grade level & \ans{Categorical} \\
Minutes of homework last night & \ans{Quantitative} \\
Rides the bus (yes / no) & \ans{Categorical} \\
Favorite lunch entr\'ee & \ans{Categorical} \\
\end{tabularx}

\vspace{4pt}
(c) The principal's first draft question is \emph{``How much homework?''}
Explain why that is not a statistical question, then rewrite it so that it is.

\ansline{It names no group and no measurement, so every student would answer it
differently.}
\ansline{``How many minutes did Bayside students spend on homework last night?''}

% ---------------------------------------------------------------- C2 (1.2)
\boxguard[16]
\item \textbf{Harbor Point Community Pool} has $1{,}500$ members. The manager
surveys $75$ members chosen at random; $27$ of them say they drive to the pool.

\vspace{2pt}
(a) Population: \ans{All 1,500 members} \quad Sample size: \ans{$75$}

\vspace{3pt}
Sample: \ans{The 75 members who were surveyed}

\vspace{3pt}
(b) What percent of the sample drives? About how many of all $1{,}500$ members
does that estimate?

\begin{work}
  \dfrac{27}{75} &= 0.36 = 36\% \\
  0.36 \times 1500 &= 540 \text{ members}
\end{work}

(c) Is the $36\%$ a parameter or a statistic? Say how you know.

\ansline{A statistic --- it describes the 75-member sample, not all 1,500 members.}

(d) Give one reason a \textbf{census} of all $1{,}500$ members is impractical here.

\ansline{Time --- reaching every one of 1,500 members would take far too long.}

% ---------------------------------------------------------------- C3 (1.3)
\boxguard[18]
\item \textbf{Riverbend High School} has $900$ students and wants a
\textbf{stratified} random sample of $60$, allocated in proportion to grade size.

\vspace{2pt}
(a) What fraction of the school is being sampled?

\begin{work}
  \dfrac{60}{900} &= \dfrac{1}{15}
\end{work}

(b) Show the work for the \textbf{10th grade} row, then complete the table.

\begin{work}
  240 \times \dfrac{1}{15} &= 16 \text{ students}
\end{work}

\vspace{-2pt}
\renewcommand{\arraystretch}{1.2}
\begin{tabularx}{\linewidth}{@{} X c c c c | c @{}}
\rowcolor{frost}
\textbf{Grade} & 9 & 10 & 11 & 12 & \textbf{Total} \\
\midrule
Students in grade & $270$ & $240$ & $210$ & $180$ & $900$ \\
Number sampled & \ans{$18$} & \ans{$16$} & \ans{$14$} & \ans{$12$} & \ans{$60$} \\
\end{tabularx}

% ---------------------------------------------------------------- C4 (1.3)
\boxguard[18]
\item \textbf{Cedar Ridge Apartments} has $600$ households on a numbered list.
The manager wants a \textbf{systematic} random sample of $60$.

\vspace{2pt}
(a) Find the interval $k$.

\begin{work}
  k &= 600 \div 60 \\
  k &= 10
\end{work}

(b) The random start is household $4$. List the first four households selected:

\vspace{2pt}
\ans{$4$} \quad \ans{$14$} \quad \ans{$24$} \quad \ans{$34$}

\vspace{4pt}
(c) Name the sampling technique each manager used.

\vspace{2pt}
\renewcommand{\arraystretch}{1.25}
\begin{tabularx}{\linewidth}{@{} X c @{}}
\rowcolor{frost}
\textbf{What the manager did} & \textbf{Technique} \\
\midrule
Drew $6$ of the $30$ buildings at random, then surveyed every
household in those $6$. & \ans{Cluster} \\
Stood at the mailboxes on Tuesday and asked whoever walked past. & \ans{Convenience} \\
Put all $600$ household numbers in a generator and drew $60$. & \ans{Simple random} \\
\end{tabularx}

% ---------------------------------------------------------------- C5 (1.4)
\boxguard[20]
\item \textbf{Millbrook Public Library} has $1{,}200$ card holders and wants to
know what fraction want Saturday evening hours. Two studies were run.

\vspace{1pt}
\textbf{Method 1 --- clipboard at the front desk.} $90$ people signed it; $72$
said yes. \\
\textbf{Method 2 --- random phone sample.} $120$ card holders called; $54$ said yes.

\vspace{2pt}
(a) Find the percent saying yes for each method.

\begin{work}
  \dfrac{72}{90} &= 0.80 = 80\% \\
  \dfrac{54}{120} &= 0.45 = 45\%
\end{work}

(b) Estimate the number of all $1{,}200$ card holders each method points to.

\begin{work}
  0.80 \times 1200 &= 960 \text{ card holders} \\
  0.45 \times 1200 &= 540 \text{ card holders}
\end{work}

(c) One method is biased. Name it, name the \textbf{type} of bias, and say which
\textbf{direction} it pushes the estimate.

\ansline{Method 1. Undercoverage --- only people already at the library on that day
could sign,}
\ansline{and they are the heaviest users, so it overestimates: $960$ against $540$.}

% ---------------------------------------------------------------- C6 (1.5)
\boxguard[20]
\item A random sample of $150$ Riverbend students is taken. Of the $60$ who use
the \textbf{Study Center}, $39$ passed all their classes. Of the $90$ who do not
use it, $72$ passed all their classes. Nobody was assigned --- students chose.

\vspace{2pt}
(a) Find each group's pass percent, and the gap between them.

\begin{work}
  \dfrac{39}{60} &= 0.65 = 65\% \\
  \dfrac{72}{90} &= 0.80 = 80\% \\
  80\% - 65\% &= 15 \text{ percentage points}
\end{work}

(b) Explanatory variable: \ans{Study Center use} \quad Response: \ans{Passed all classes}

\vspace{3pt}
(c) Using all $150$ students, find the overall pass percent and estimate how
many of all $900$ students that is.

\begin{work}
  \dfrac{39 + 72}{150} &= \dfrac{111}{150} = 0.74 = 74\% \\
  0.74 \times 900 &= 666 \text{ students}
\end{work}

(d) Name one \textbf{lurking variable} that could produce this result, and say
which way it would push the numbers.

\ansline{Students who were already struggling are the ones who go to the Study
Center.}
\ansline{That pushes the users' pass rate down without the Study Center causing it.}

% ---------------------------------------------------------------- C7 (1.6)
\boxguard[20]
\item Riverbend now runs an \textbf{experiment}. Of $120$ volunteers, a random
number generator sends $60$ to required weekly Study Center sessions and $60$ to
no sessions. Afterwards $45$ of the session group and $33$ of the other group
passed all their classes.

\vspace{2pt}
(a) Find each group's pass percent and the gap.

\begin{work}
  \dfrac{45}{60} &= 0.75 = 75\% \\
  \dfrac{33}{60} &= 0.55 = 55\% \\
  75\% - 55\% &= 20 \text{ percentage points}
\end{work}

(b) Treatment: \ans{Required weekly sessions} \quad Control group: \ans{The 60 with no sessions}

\vspace{3pt}
(c) Instead, Riverbend \textbf{blocks} on last quarter's grades: $80$ volunteers
passed everything and $40$ did not. Randomly split each block in half. How many
go to each group?

\begin{work}
  80 \div 2 &= 40 \text{ per group (passed-everything block)} \\
  40 \div 2 &= 20 \text{ per group (did-not block)}
\end{work}

% ---------------------------------------------------------------- C8 (1.8)
\boxguard[30]
\item The Riverbend Student Council asked its stratified sample of $60$ students
how to spend a \$$3{,}000$ grant. Complete the table, then complete the bar graph.

\vspace{2pt}
\renewcommand{\arraystretch}{1.25}
\begin{tabularx}{\linewidth}{@{} X c c @{}}
\rowcolor{frost}
\textbf{Choice} & \textbf{Count} & \textbf{Relative frequency} \\
\midrule
Club funding    & $21$ & \ans{$35\%$} \\
Field trips     & $18$ & \ans{$30\%$} \\
Gym equipment   & $15$ & \ans{$25\%$} \\
Longer lunch    & $6$  & \ans{$10\%$} \\
\end{tabularx}

\vspace{6pt}
\begin{center}
\begin{tikzpicture}[x=1cm, y=0.075cm]
  \foreach \y in {5,15,25,35}
    \draw[linegray, very thin] (0,\y) -- (10.4,\y);
  \foreach \y in {10,20,30,40}
    \draw[linegray] (0,\y) -- (10.4,\y);
  \foreach \y in {0,10,20,30,40}
    \node[left, font=\scriptsize] at (-0.05,\y) {\y\%};
  \foreach \x in {2.6,5.2,7.8}
    \draw[linegray, dashed] (\x,0) -- (\x,45);
  \foreach \i/\lab in {1/Club funding, 2/Field trips, 3/Gym equipment,
                       4/Longer lunch}
    \node[below, font=\scriptsize, align=center, text width=2.4cm]
      at ({2.6*\i-1.3},-0.5) {\lab};
  \draw[cerulean, thick] (0,45) -- (0,0) -- (10.4,0);
  % Key only: the four correct bars, drawn inside the existing bounding box.
  \foreach \i/\h in {1/35, 2/30, 3/25, 4/10}
    \fill[keyred, opacity=0.55] ({2.6*\i-2.0},0) rectangle ({2.6*\i-0.6},\h);
\end{tikzpicture}
\end{center}

\vspace{2pt}
About how many of all $900$ Riverbend students would choose club funding?

\begin{work}
  0.35 \times 900 &= 315 \text{ students}
\end{work}

\end{enumerate}
\normalsize

% ======================================================================
\boxguard[10]
\parthead{Part D --- Extended Response \hfill 3 questions, 10 points each}

\small
\begin{enumerate}[leftmargin=2.2em, itemsep=10pt, topsep=2pt, start=31]

% ---------------------------------------------------------------- D1
\item Riverbend has now run \textbf{two} studies of the same Study Center.

\vspace{1pt}
\textbf{Study A.} A random sample of students was watched. Of the students who
\emph{chose} to use the Study Center, $65\%$ passed all classes; of those who did
not, $80\%$ did. \\
\textbf{Study B.} Of $120$ volunteers, half were \emph{assigned} to weekly
sessions. Of those, $75\%$ passed all classes; of the others, $55\%$ did.

\vspace{3pt}
(a) Which study is observational and which is an experiment? Give the one
feature that decides it.

\ansline{A is observational, B is an experiment. The deciding feature is who decides
the groups:}
\ansline{in A the students chose; in B a random generator assigned them.}

(b) Study A's gap is $15$ points and Study B's is $20$ points. Which study lets
the council write \emph{``the Study Center caused more students to pass''}?
Justify your choice.

\ansline{Study B. Random assignment makes the two groups alike at the start, so the
sessions are}
\ansline{the only difference left. Study A's larger-looking gap earns only ``is
associated with.''}

(c) Explain how the same Study Center can look \emph{harmful} in Study A and
\emph{helpful} in Study B.

\ansline{In A the students who chose the Study Center were already struggling, so
their group}
\ansline{started behind. In B nobody chose, so that head start cannot be hiding in
the gap.}

% ---------------------------------------------------------------- D2
\boxguard[26]
\item \textbf{Cedar Ridge Apartments} mailed a survey about a proposed dog park
to all $600$ households. $150$ were returned, and $96$ of those said yes. The
management office's own sign-up records show that $240$ of the $600$ households
actually want one.

\vspace{2pt}
(a) Find the response rate, the survey's percent saying yes, and the number of
households that percent points to.

\begin{work}
  \dfrac{150}{600} &= 0.25 = 25\% \text{ response rate} \\
  \dfrac{96}{150} &= 0.64 = 64\% \\
  0.64 \times 600 &= 384 \text{ households}
\end{work}

(b) The records say the truth is $240$ of $600$, or $40\%$. Name the bias, and
give its size in percentage points \emph{and} in households.

\ansline{Nonresponse bias --- the $75\%$ who threw the survey away cared less about a
dog park.}
\ansline{The survey overestimates by $24$ percentage points, or $384 - 240 = 144$
households.}

(c) The manager mails again and gets $300$ returns, $192$ of them yes. Compute
that percent and explain what the result shows.

\begin{work}
  \dfrac{192}{300} &= 0.64 = 64\%
\end{work}

\ansline{Twice the responses, the same wrong $64\%$: sample size does not remove bias.}

(d) Describe one change that would actually reduce the bias, and name the bias
it removes.

\ansline{Draw a random sample of households and follow up by phone or in person until
nearly}
\ansline{all of them answer. That removes the nonresponse bias, not the sample size.}

% ---------------------------------------------------------------- D3
\boxguard[24]
\item \textbf{Bayside Middle School} ($750$ students) wants to know whether to
add a late activity bus. The office has a numbered list of all $750$ students,
and the four grades are believed to differ a lot in their answers.

\vspace{2pt}
(a) Choose \textbf{one} of the four sampling techniques and justify it from what
this context allows.

\ansline{Stratified. A full numbered list exists, and the grades are expected to
differ, so}
\ansline{sampling within each grade guarantees all four are represented in
proportion.}

(b) Name one of the five \textbf{data collection methods} and say why it fits
this question.

\ansline{Survey. The question has a small set of fixed answers and needs many students,}
\ansline{so a short written instrument can be counted quickly.}

(c) Instead, a teacher asks her own homeroom of $30$ students; $12$ say yes. She
reports: \emph{``$12 \div 30 = 40\%$, and $0.40 \times 750 = 300$ students want
the late bus.''} Her arithmetic is correct. Explain what is wrong with the
report, and state what she \emph{is} allowed to say.

\ansline{Her homeroom is a convenience sample --- chance never chose it, and the other
$720$}
\ansline{students had no possibility of being asked, so the $300$ is not an estimate
for the school.}
\ansline{She may say only that $40\%$ of her own homeroom wants the late bus.}

\end{enumerate}
\normalsize

\end{document}
